\chapter{Chaos Diagnostics and Model Validation}

\section{Model‑Implied vs. Empirical Diagnostics}
Using the estimated parameters, we simulate a long synthetic path ($50\,000$ observations) and recompute the tail index, Hurst exponent, and multifractal width. Table \ref{tab:diag} compares these model‑implied values with the empirical ones (post‑break sample).

\begin{table}[H]
\centering
\caption{Empirical vs. model‑implied chaos diagnostics (post‑break sample).}
\label{tab:diag}
\begin{tabular}{l c c c}
\toprule
\textbf{Diagnostic} & \textbf{Empirical} & \textbf{Model‑implied} & \textbf{Match} \\
\midrule
Tail index $\gamma$ (Hill, $|r_t|$) & $2.7805$ & $2.8928$ & Yes \\
Hurst exponent $H$ (DFA on $|r_t|$) & $0.7751$ & $0.8260$ & Yes \\
Multifractal width $\Delta\alpha$ & $0.7959$ & $0.8096$ & Yes \\
\bottomrule
\end{tabular}
\end{table}

The model captures the empirical tail heaviness, long memory, and multifractal structure remarkably well. The slight overestimation of $H$ and $\Delta\alpha$ is within sampling variability.

\section{Conditions for Chaotic Regime}
Recall the theoretical conditions for the chaotic phase:
\begin{enumerate}[label=(\roman*)]
    \item Near‑critical branching: $B \in [0.9,1)$.
    \item Heavy volatility tails: $\gamma \le 2$.
    \item Long memory: $H > 0.5$.
    \item Multifractality: $\Delta\alpha > 0.1$.
\end{enumerate}
Our results show:
\begin{itemize}
    \item $B = 0.7468$ (not satisfied);
    \item $\gamma = 2.7805 > 2$ (not satisfied);
    \item $H = 0.7751 > 0.5$ (satisfied);
    \item $\Delta\alpha = 0.7959 > 0.1$ (satisfied).
\end{itemize}
Thus only two of the four conditions are met. The NEPSE banking sector does not exhibit the full chaotic regime; however, it displays strong long memory and multifractality, which are already sufficient to justify the COGARCH framework for risk management. The branching ratio being below $0.9$ suggests that the self‑exciting feedback is weaker than in mature markets, possibly due to lower liquidity and fewer high‑frequency traders.

\section{Filtered Volatility and Jump Detection}
Using the GARCH(1,1) approximation of the COGARCH,
\[
\omega = \kappa\theta = 1.5567\times10^{-5},\quad
\alpha_1 = \phi\lambda = 0.0545,\quad
\beta_1 = 1-\kappa = 0.9271,
\]
the persistence is $\alpha_1+\beta_1 = 0.9815$, confirming strong volatility clustering. Applying a $3\sigma$ threshold to the standardised residuals identifies jump days; the time‑varying conditional variance yields a higher detection rate than the simple rolling filter.

